% !LW recipe = make-deck
% Level 2F (programs track, FPIL): selection. Theory: conditions,
% selection, select on lists. Lean: build select; law 1 on an example.
% Build: make (rules in ../common/deck.mk).
\DocumentMetadata{
  lang        = en-GB,
  pdfstandard = ua-2,
  tagging     = on,
}
\documentclass{ltx-talk}

\input{preamble}

\title{Level 2F: Selection, as a Program}
\author{The Proofs and Programs Club}
\institute{Royal Holloway, University of London}

\begin{document}

\leveltitle{2F}{Selection, as a Program}{Programs track: follows FPIL}
  {After \levellink{1}}{Next: \levellink{3f}}

\levelframe{2F}{Level 2F at a Glance}{%
  \code{select}, which keeps the rows of a table that satisfy a
  condition.}{%
  facts about selections of the club's tables, by computing.}{%
  define a function by cases on a list: empty, or a first row and the rest
    & \rds{fpil-data} \\
  write one function for tables of any row type & \rds{fpil-poly} \\
  pass a function, such as a condition, as an argument & \rds{fpil-fun} \\
}

% ---------------------------------------------------------------- theory
\theorypart{Conditions, selection, and selection on lists}

\begin{frame}
  \frametitle{A Condition Says Yes or No to Each Row}
  Write $\B = \{\true, \false\}$ for the two truth values.
  \defn{Definition 2F.1 (Condition)}{A \emph{condition} on rows of
    $D_1 \times D_2$ is a function $p : D_1 \times D_2 \to \B$.}
  For example, on $\rel{Member}$ and on $\rel{Attended}$:
  \[
    \mathit{cs}(n, d) = \true \iff d = \text{CS}, \qquad
    \mathit{late}(n, w) = \true \iff w \geq 2.
  \]
  \defn{Definition 2F.2 (Both conditions)}{$(p \wedge q)(t) = \true$
    exactly when $p(t) = \true$ and $q(t) = \true$.}
  \theorytag{2F}
\end{frame}

\begin{frame}
  \frametitle{Selection Keeps the Rows That Satisfy a Condition}
  \defn{Definition 2F.3 (Selection)}{For a table $R$ and a condition $p$,
    \ $\sel{p}(R) = \{t \in R \mid p(t) = \true\}$.}
  {\centering
    \smalltable{$\sel{\mathit{cs}}(\rel{Member})$}{ll}
      {$\att{name}$ & $\att{dept}$}
      {Alan & CS \\ Grace & CS}\hspace{16mm}
    \smalltable{$\sel{\mathit{late}}(\rel{Attended})$}{lr}
      {$\att{name}$ & $\att{week}$}
      {Alan & 2 \\ Grace & 3}\par}
  \medskip
  A selection never adds rows: $\sel{p}(R) \subseteq R$. In SQL, a
  selection is a \code{WHERE} clause.
  \slideref{\lnk{ahv4}}
  \theorytag{2F}
\end{frame}

\begin{frame}
  \frametitle{On a List, Selection Is Three Equations}
  \defn{Definition 2F.4 (Selection on lists)}{For a condition $p$:
    \[
      \begin{array}{@{}l@{\;=\;}l@{\qquad}l@{}}
        \mathit{select}\ p\ [\,] & [\,] \\
        \mathit{select}\ p\ (t :: l) & t :: \mathit{select}\ p\ l &
          \text{if } p(t) = \true \\
        \mathit{select}\ p\ (t :: l) & \mathit{select}\ p\ l &
          \text{if } p(t) = \false \\
      \end{array}
    \]}
  Here $t :: l$ is the list with first row $t$, then the list $l$. For
  example:
  \[
    \mathit{select}\ \mathit{late}\ [(\text{Alan}, 1), (\text{Alan}, 2)]
    = \mathit{select}\ \mathit{late}\ [(\text{Alan}, 2)]
    = [(\text{Alan}, 2)].
  \]
  \defn{Fact 2F.5}{If $l$ represents $R$ (Definition 1.2), then
    $\mathit{select}\ p\ l$ represents $\sel{p}(R)$. Level 6 proves it.}
  \theorytag{2F}
\end{frame}

% ---------------------------------------------------------------- Lean
\leanpart{Build select, and run it}

\begin{frame*}
  \frametitle{Your Turn: Build select}
  \leanblock{select.lean}
  \begin{itemize}
    \item A condition is a function to \code{Bool}, Lean's $\B$:
      \code{fun row => row.2 == 1}.
    \item \ex{2F.1} Replace the \code{sorry}s with Definition 2F.4: use
      \code{if p row then \ldots{} else \ldots}.
    \item \ex{2F.2} Prove that \code{select} keeps 2 rows of week 1, by
      \code{rfl}. \ex{2F.3} \code{\#eval} the CS members.
  \end{itemize}
  \slideref{\lnk{fpil-data} \rfsep \lnk{fpil-poly}}
  \leantag{2F}
\end{frame*}

\begin{frame*}
  \frametitle{Law 1, Checked by Lean on an Example}
  \leanblock{law.lean}
  \begin{itemize}
    \item \textbf{Law 1}: $\sel{p}(\sel{q}(R)) = \sel{p \wedge q}(R)$:
      two selections are one. Level 2T proves it for every table; here
      Lean checks one case, by computing.
    \item \ex{2F.4} Check law 1 on another example of your own.
    \item \ex{2F.5} Is $\sel{p}(\sel{q}(R)) = \sel{q}(\sel{p}(R))$?
      Check it in Lean on an example.
  \end{itemize}
  \leantag{2F}
\end{frame*}

\end{document}
